\documentclass[11pt]{amsart}
\usepackage[margin=1.1in]{geometry}
\usepackage{amsmath,amssymb,amsthm,booktabs,array,url,enumitem}
\usepackage{algorithm,algpseudocode}
\usepackage[dvipsnames]{xcolor}
\usepackage[pdftitle={Exhaustive searches for Williams numbers: the four-prime case and a general 10^24 bound},pdfauthor={Raz Vakil},pdfkeywords={Carmichael numbers, Lucas-Carmichael numbers, Williams numbers, Baillie-PSW, exhaustive search}]{hyperref}
\hypersetup{colorlinks=true,linkcolor=blue!50!black,citecolor=blue!50!black,urlcolor=blue!50!black}
\newtheorem{theorem}{Theorem}[section]
\newtheorem{lemma}[theorem]{Lemma}
\newtheorem{proposition}[theorem]{Proposition}
\newtheorem{corollary}[theorem]{Corollary}
\theoremstyle{definition}
\newtheorem{definition}[theorem]{Definition}
\theoremstyle{remark}
\newtheorem{remark}[theorem]{Remark}
\providecommand{\lcm}{\operatorname{lcm}}
\DeclareMathOperator{\omg}{\omega}

\title[Exhaustive searches for Williams numbers]{Exhaustive searches for Williams numbers:\\ the four-prime case and a general $10^{24}$ bound}
\author{Raz Vakil}
\address{Tel Aviv, Israel}
\email{Raz.Vakil@gmail.com}
\subjclass[2020]{Primary 11Y11; Secondary 11A51, 11Y16}
\keywords{Carmichael numbers, Lucas--Carmichael numbers, Williams numbers, Baillie--PSW test, exhaustive search}
\date{1 November 2026}

\begin{document}

\begin{abstract}
An odd squarefree composite integer is a Williams number if it is both Carmichael and Lucas--Carmichael; none is known. For a four-prime Williams number $n=p_1p_2p_3p_4$ we give a self-contained restatement, with a slight sharpening, of the bounds due to McIntosh and Dipra, and derive from them an enumeration that fixes $(p_1,p_2)$, places $p_3p_4$ in a single residue class modulo $\lcm\bigl((p_1^2-1)/2,(p_2^2-1)/2\bigr)$, and rejects almost every candidate with one multiplication. Certified computations prove that no such number satisfies $n\le10^{30}$ or $p_1\le500{,}000$, extending McIntosh and Dipra's $n>10^{24}$, $p_1>2\times10^4$. Separately, a factor-and-test scan of all $308{,}279{,}939$ Carmichael numbers below $10^{24}$ finds no Williams number. Code, logs and certification records are archived.
\end{abstract}

\maketitle

\section{Introduction}

Throughout, $n$ denotes an odd squarefree composite integer with prime factorisation $n=p_1p_2\cdots p_k$, $p_1<\dots<p_k$, $k=\omg(n)\ge2$.

\begin{definition}
$n$ is a \emph{Carmichael number} if $(p-1)\mid(n-1)$ for every $p\mid n$ (Korselt's criterion), and a \emph{Lucas--Carmichael number} if $(p+1)\mid(n+1)$ for every $p\mid n$. A number that is both is a \emph{Williams number} (a $1$-Williams number in Echi's terminology \cite{Echi07}).
\end{definition}

Carmichael numbers are the absolute Fermat pseudoprimes; Lucas--Carmichael numbers play the same role for Lucas sequences with discriminant a quadratic non-residue of every $p\mid n$. H.~C.~Williams \cite{Williams77} asked in 1977 whether a number can be both, proved that it cannot have exactly three prime factors and cannot be divisible by $3$, and showed that under his Lucas-sequence formulation the number of prime factors must be odd (in the divisibility formulation above, all $p_i$ lie in one class modulo $4$, and $k$ is odd when that class is $3$; Lemma~\ref{lem:classes}). Echi \cite{Echi07} named such numbers Williams numbers, within a family of $a$-Williams numbers, and dropped the parity restriction; Bouall\`egue, Echi and Pinch \cite{BEP10} and Al-Rasasi and Ghanmi \cite{AG13} studied general $a$; Pinch \cite{Pinch07} reported that there is no Williams number up to $10^{16}$ with an odd number of prime factors. McIntosh \cite{McIntosh14} treated the parallel problem $(p\pm1)\mid(n-1)$, and McIntosh and Dipra \cite{MD15} the Williams numbers themselves. Zhang \cite{Zhang15} developed counting heuristics for Williams numbers built from a fixed set of primes and offered a prize for the first example. No Williams number is known.

The question is tied to primality testing. Pomerance's 1984 heuristic argument that counterexamples to the Baillie--PSW primality test exist \cite{Pomerance84}, restated in \cite{BFW21}, builds $n=\prod p$ from primes with $p\equiv3\pmod 8$, $(5/p)=-1$ and smooth squarefree $(p-1)/2$, $(p+1)/4$, chosen so that $(p-1)\mid(n-1)$ and $(p+1)\mid(n+1)$ for all $p\mid n$; such an $n$ is in particular a Williams number \cite{McIntosh14,MD15,Zhang15}. The construction attempted by Chen and Greene \cite{ChenGreene03} for Grantham's \$620 problem \cite{Grantham} uses weaker order conditions; the related set of $2{,}030$ primes does satisfy $(p-1)\mid M$, $(p+1)\mid N$, and an unrefereed report by Phatak \cite{Phatak20} states that no subset of that set is a Williams number. The Lenstra--Pomerance analysis of Agrawal's conjecture from the AKS primality test \cite{AKS,LP} shows that the natural counterexamples are Williams numbers whose primes are all $\equiv3\pmod{80}$. Howe's higher-order Carmichael numbers \cite{Howe00} are a further neighbour: by Howe's criterion a squarefree $n$ is Carmichael of order $2$ iff for every $p\mid n$ one has $n\equiv1$ or $n\equiv p\pmod{p^2-1}$, while $n$ is a Williams number iff $n\equiv p\pmod{(p^2-1)/2}$ for every $p\mid n$ (Lemma~\ref{lem:quotient}).

\paragraph{Prior results for four prime factors, and what is new here.} McIntosh and Dipra \cite{MD15} proved that for a fixed set of $d-3$ prime factors only finitely many Williams numbers with $d$ prime factors exist, with explicit bounds on the remaining three primes, observed that all prime factors are congruent modulo $12$ (and $\equiv1$ when their number is even), organised two searches, and showed that a Williams number with exactly four prime factors has $p_1>2\times10^4$ and $n>10^{24}$. Table~\ref{tab:compare} sets their results for $d=4$ beside ours. Theorem~\ref{thm:A} is their theorem restated with a self-contained proof; the new contributions are the enumeration algorithm of \S\ref{sec:alg}, its certified completeness (\S\ref{sec:complete}), the resulting exclusions (Theorem~\ref{thm:B}), and the general $10^{24}$ census (Theorem~\ref{thm:bound}).

\begin{table}[h]
\centering\small
\begin{tabular}{p{2.6cm}p{5.6cm}p{5.6cm}}
\toprule
 & McIntosh--Dipra \cite{MD15} & this paper\\
\midrule
all $p_i\equiv1\ (12)$ & \S2 & Lemma~\ref{lem:classes}(c)\\
$p_2$ & $p_2<5p_1^{3/2}$ (Thm.~3) & $p_2^2\le24p_1^3+1$ (Thm.~\ref{thm:A})\\
$p_3$ & $p_3<p_1p_2$ (Thm.~3) & $p_3<2p_1p_2$ (weaker; sufficient)\\
$p_4$ & $p_4<\sqrt{2p_1p_2p_3}$ (Thm.~3) & same\\
search method & loops over the three quotients; Pinch-style loops over compatible primes & CRT class for $p_3p_4$ and window filter (\S\ref{sec:alg}); certified completeness\\
four primes excluded & $p_1\le2\times10^4$; $n\le10^{24}$ & $p_1\le500{,}000$; $n\le10^{30}$ (Thm.~\ref{thm:B})\\
all $d$ & searches for $4\le d\le13$ at unstated bounds in $[10^{20},10^{35}]$ & none below $10^{24}$, any $d$ (Thm.~\ref{thm:bound})\\
\bottomrule
\end{tabular}
\caption{Four-prime Williams numbers: prior and present results.}\label{tab:compare}
\end{table}

All code, logs and certification records are available (\S\ref{sec:repro}). Appendix~\ref{sec:obstr} proves that Chernick-type polynomial families whose divisibilities hold identically in the parameter cannot produce Williams numbers with three factors, with uniform offsets $\pm1$, or in generic position; a supplementary note \cite{Supp} records two heuristic experiments.

\section{Preliminaries}\label{sec:prelim}

For an odd prime $p$ write $L_p=(p^2-1)/2=\lcm(p-1,p+1)$.

\begin{lemma}\label{lem:quotient}
Let $p\mid n$. Then $(p-1)\mid(n-1)$ and $(p+1)\mid(n+1)$ hold together if and only if $n/p\equiv1\pmod{L_p}$, equivalently $n\equiv p\pmod{L_p}$.
\end{lemma}
\begin{proof}
The two divisibilities say $n\equiv1\pmod{p-1}$ and $n\equiv-1\pmod{p+1}$. Since $p\equiv1\pmod{p-1}$ and $p\equiv-1\pmod{p+1}$, this is $n\equiv p$ modulo $\lcm(p-1,p+1)=L_p$ (the gcd of $p\mp1$ is $2$). Finally $L_p\mid n-p=p\,(n/p-1)$ with $\gcd(p,L_p)=1$ gives $L_p\mid n/p-1$, and conversely.
\end{proof}

\begin{lemma}[congruence classes]\label{lem:classes}
Let $n$ be a Williams number with $k$ prime factors.
\begin{enumerate}[label=\textup{(\alph*)}]
\item All $p\mid n$ lie in one class $c$ modulo $4$, and $n\equiv c\pmod 4$; hence $k$ is odd if $c=3$.
\item $3\nmid n$; all $p\mid n$ lie in one class $c'$ modulo $3$, and $n\equiv c'\pmod3$; hence $k$ is odd if $c'=2$.
\item If $k$ is even, every $p\mid n$ is $\equiv1\pmod{12}$.
\end{enumerate}
\end{lemma}
\begin{proof}
(a) $8\mid p^2-1$ for odd $p$, so $4\mid L_p$ and $n\equiv p\pmod 4$ for every $p\mid n$ by Lemma~\ref{lem:quotient}. Thus all $p\equiv c$, and $n\equiv c^k\pmod4$; if $c=3$ this forces $k$ odd. (b) For $p\ne3$, $3\mid p^2-1$, so $3\mid L_p$ and $n\equiv p\pmod3$. If $3\mid n$ then every other prime factor would be $\equiv n\equiv0\pmod3$, impossible for $k\ge2$; so $3\nmid n$, all $p\equiv c'\pmod3$ and $n\equiv c'^k\equiv c'$, which for $c'=2$ forces $k$ odd. (c) is (a) and (b) with $k$ even.
\end{proof}

Part (c) is observed in \cite[\S2]{MD15}. Williams proved that $k=3$ is impossible; the following proof is essentially that of \cite[Thm.~1]{MD15}.

\begin{proposition}[Williams]\label{prop:k3}
No Williams number has exactly three prime factors.
\end{proposition}
\begin{proof}
Let $n=p_1p_2p_3$. By Lemma~\ref{lem:quotient} at $p_3$, $p_1p_2-1=tL_3$ with $t\ge1$ (as $p_1p_2>1$); since $p_1p_2<p_3^2$ we have $t<2$, so $2p_1p_2=p_3^2+1$. At $p_2$, $p_1p_3-1=sL_2$ with $s\ge1$; now $p_3<\sqrt{2p_1p_2}<\sqrt2\,p_2$, so $p_1p_3<\sqrt2\,p_2^2$ and $s<2\sqrt2\,p_2^2/(p_2^2-1)<3$ (as $p_2\ge5$). If $s=2$ then $p_1p_3=p_2^2$, impossible; so $2p_1p_3=p_2^2+1$. Subtracting, $2p_1(p_2-p_3)=p_3^2-p_2^2$, i.e. $2p_1=-(p_2+p_3)<0$, a contradiction.
\end{proof}

\section{Four prime factors: structural bounds}\label{sec:k4}

\begin{theorem}[McIntosh--Dipra \cite{MD15}, $d=4$; restated]\label{thm:A}
Let $n=p_1p_2p_3p_4$, $p_1<p_2<p_3<p_4$, be a Williams number, and write $P=p_1p_2$, $q=p_3$, $r=p_4$.
\begin{enumerate}[label=\textup{(\alph*)}]
\item Every $p_i\equiv1\pmod{12}$; in particular $p_1\ge13$ and $p_2\ge37$.
\item There are integers $d,e,s\ge1$ with
\begin{equation}\label{eq:three}
2Pq=d(r^2-1)+2,\qquad 2Pr=e(q^2-1)+2,\qquad 2p_1qr=s(p_2^2-1)+2,
\end{equation}
and consequently $q<2P$, $r<\sqrt{2Pq}<2P$, $m:=qr<4P^2$, $n<4P^3$ and $d<2p_1$.
\item $p_2^2\le24p_1^3+1$.
\end{enumerate}
Consequently $p_2<4.9\,p_1^{3/2}$, $p_3,p_4<2P<9.8\,p_1^{5/2}$ and $n<477\,p_1^{15/2}$.
\end{theorem}

\begin{proof}
(a) is Lemma~\ref{lem:classes}(c).

(b) Lemma~\ref{lem:quotient} at $r$ gives $n/r-1=Pq-1=dL_r$ with an integer $d\ge1$, which is the first equation; the second and third are Lemma~\ref{lem:quotient} at $q$ and at $p_2$. From $2Pq=d(r^2-1)+2\ge r^2+1>q^2$ we get $q<2P$ and $r^2<2Pq$, hence $r<\sqrt{2Pq}<2P$, $m=qr<4P^2$ and $n=Pm<4P^3$. Also $d=2(Pq-1)/(r^2-1)<2Pq/(r^2-1)<2P/q<2p_1$ because $r\ge q+1$ and $q>p_2$.

(c) Multiply the three equations in~\eqref{eq:three}. The left side is $8P^2p_1\,q^2r^2=8p_1^3p_2^2q^2r^2$. Writing $d(r^2-1)+2=dr^2\bigl(1-\tfrac{d-2}{dr^2}\bigr)$ and similarly for the other two factors, the right side is $d\,e\,s\;r^2q^2p_2^2\,\Pi$ with
\[
\Pi=\Bigl(1-\frac{d-2}{dr^2}\Bigr)\Bigl(1-\frac{e-2}{eq^2}\Bigr)\Bigl(1-\frac{s-2}{sp_2^2}\Bigr).
\]
Hence
\begin{equation}\label{eq:des}
d\,e\,s\cdot\Pi=8p_1^3 .
\end{equation}
For an integer $t\ge1$ and real $x>0$ one has $-1/x^2\le (t-2)/(tx^2)<1/x^2$, so each factor of $\Pi$ lies in $[1-1/x^2,\,1+1/x^2]$ with $x\in\{r,q,p_2\}$, $x\ge p_2$; therefore
\begin{equation}\label{eq:Pi}
(1-p_2^{-2})^3\le\Pi\le(1+p_2^{-2})^3 .
\end{equation}
Now reduce the three equations~\eqref{eq:three} modulo $p_1$. Each left side is divisible by $p_1$, so $d(r^2-1)\equiv e(q^2-1)\equiv s(p_2^2-1)\equiv-2\pmod{p_1}$. As $p_1$ is odd, $p_1\nmid d$, $p_1\nmid e$ and $p_1\nmid s$. Hence $des\ne8p_1^3$, i.e. $\Pi\ne1$.

\emph{Case $\Pi<1$.} Then $des\ge8p_1^3+1$ and $\Pi=8p_1^3/des\le1-\frac1{8p_1^3+1}$; with $\Pi\ge(1-p_2^{-2})^3\ge1-3p_2^{-2}$ this gives $p_2^2\le3(8p_1^3+1)=24p_1^3+3$. Since $p_2\equiv1\pmod{12}$ implies $p_2^2\equiv1\pmod{24}$, while $24p_1^3+2$ and $24p_1^3+3$ are $\equiv2,3\pmod{24}$, in fact $p_2^2\le24p_1^3+1$.

\emph{Case $\Pi>1$.} A factor $1-\frac{t-2}{tx^2}$ of $\Pi$ exceeds $1$ only when $t=1$. Now $e>d\ge1$, because $e/d=\frac{(Pr-1)(r^2-1)}{(Pq-1)(q^2-1)}>1$ for $r>q$; and $s>2p_1$, because $qr\ge p_2^2+1$ gives $s=\frac{2(p_1qr-1)}{p_2^2-1}\ge\frac{2(p_1p_2^2+p_1-1)}{p_2^2-1}>2p_1$. So only the first factor can exceed $1$, and $\Pi\le1+r^{-2}$. On the other hand $des\le8p_1^3-1$ gives $\Pi=8p_1^3/des\ge1+\frac1{8p_1^3-1}$. Hence $r^2\le8p_1^3-1$, and $p_2^2<r^2<8p_1^3\le24p_1^3+1$.

In both cases $p_2^2\le24p_1^3+1$. The consequences follow from (b) and (c): $p_2^2\le24p_1^3+1<24.01\,p_1^3=(4.9\,p_1^{3/2})^2$ for $p_1\ge5$, so $p_2<4.9\,p_1^{3/2}$; then $q,r<2P<9.8\,p_1^{5/2}$ and $n<4P^3<4\cdot4.9^3\,p_1^{15/2}<477\,p_1^{15/2}$.
\end{proof}

\begin{remark}
In the notation of \cite{MD15} ($P,q,r,s$ and $f_s,f_r,f_q$ for our $p_1,p_2,p_3,p_4$ and $d,e,s$), Theorems~2 and~3 there give $q<5P^{3/2}$, $r<Pq$ and $s<\sqrt{2Pqr}$, and show by Newton's inequality that the quotient product always exceeds $8P^3$, so that the case $\Pi>1$ above never occurs. The statement above differs only in the slightly sharper bound on $p_2$ (from the congruence modulo $24$) and the weaker bound on $p_3$, which is all the enumeration needs.
\end{remark}

\begin{corollary}\label{cor:A}
For each prime $p_1$ there are only finitely many four-prime Williams numbers with smallest prime $p_1$, all below $477\,p_1^{15/2}$. Conversely a four-prime Williams number $n$ has $p_1>(n/477)^{2/15}$.
\end{corollary}

\section{The enumeration algorithm}\label{sec:alg}

Theorem~\ref{thm:A} makes two kinds of complete enumeration possible: all $n\le B$ (where $n>p_1p_2^3$ forces $p_1<B^{1/4}$ and $p_2<(B/p_1)^{1/3}$), and all $n$ whatsoever with $p_1\le X$ (where $p_2\le\sqrt{24p_1^3+1}$). Both reduce to a search over pairs $(p_1,p_2)$ of primes $\equiv1\pmod{12}$.

\begin{lemma}[one class for the cofactor]\label{lem:crt}
With notation as in Theorem~\ref{thm:A}, $m=qr$ satisfies $p_2m\equiv1\pmod{L_{p_1}}$ and $p_1m\equiv1\pmod{L_{p_2}}$. These are solvable iff $p_1\nmid p_2^2-1$ and $p_1\equiv p_2\pmod{\gcd(L_{p_1},L_{p_2})}$, in which case $m$ lies in a single residue class modulo $\Lambda=\lcm(L_{p_1},L_{p_2})$.
\end{lemma}
\begin{proof}
The congruences are Lemma~\ref{lem:quotient} at $p_1$ and $p_2$ ($n/p_1=p_2m$, $n/p_2=p_1m$). The first is always solvable: $p_2\mid L_{p_1}$ would need $p_2\mid p_1\pm1$ with $p_2>p_1$, impossible for odd primes. The second is solvable iff $\gcd(p_1,L_{p_2})=1$, i.e. $p_1\nmid p_2^2-1$. Given both, let $g=\gcd(L_{p_1},L_{p_2})$; the pair is compatible iff $p_2^{-1}\equiv p_1^{-1}\pmod g$, i.e. $p_1\equiv p_2\pmod g$ (note $\gcd(p_1p_2,g)=1$). The rest is the Chinese remainder theorem.
\end{proof}

Since $\Lambda\approx P^2/(4g)$ while $m<4P^2$, Lemma~\ref{lem:crt} leaves $O(g)$ candidates $m$ per pair. The next lemma rejects almost all of them with one multiplication.

\begin{lemma}[window]\label{lem:window}
With $d,e$ as in~\eqref{eq:three} and $j:=de$, $T:=4P^2$,
\begin{equation}\label{eq:window}
|jm-T|<\max\bigl(2j\sqrt{2p_1},\;8p_1\bigr).
\end{equation}
Moreover, for fixed $m>2\lceil\sqrt{2p_1}\rceil$ (which holds for $m>p_2^2$) the set of real $j>0$ satisfying $|jm-T|\le\max(2j\lceil\sqrt{2p_1}\rceil,8p_1)$ is an interval containing $T/m$; hence it contains an integer iff it contains $\lfloor T/m\rfloor$ or $\lceil T/m\rceil$.
\end{lemma}
\begin{proof}
From~\eqref{eq:three}, $d=2(Pq-1)/(r^2-1)$ and $e=2(Pr-1)/(q^2-1)$, so $j\,(q^2-1)(r^2-1)=4(Pq-1)(Pr-1)$. Expanding with $m=qr$,
\[
jm^2-Tm=j(q^2+r^2-1)-4P(q+r)+4 .
\]
The right side is a difference of two positive numbers, so its absolute value is at most $\max\bigl(j(q^2+r^2-1),\,4P(q+r)-4\bigr)\le\max(2jr^2,\,8Pr)$. Dividing by $m=qr$ gives $|jm-T|\le\max(2jr/q,\,8P/q)$. By Theorem~\ref{thm:A}(b), $r^2<2Pq$, so $r/q<\sqrt{2P/q}<\sqrt{2p_1}$ (as $q>p_2$), and $P/q<p_1$; this is~\eqref{eq:window}. For the interval property, put $S=\lceil\sqrt{2p_1}\rceil$ and $\phi(j)=|jm-T|-\max(2jS,8p_1)$. On $j\ge T/m$ the first term increases with slope $m$ and the second with slope at most $2S<m$, so $\phi$ is strictly increasing; on $j\le T/m$ the first term decreases and the second is non-decreasing, so $\phi$ is strictly decreasing. Thus $\{\phi\le0\}$ is an interval, and it contains $T/m$ since $\phi(T/m)<0$.
\end{proof}

\begin{lemma}[reconstruction]\label{lem:cbrt}
Let $n=p_1p_2qr$ be a four-prime Williams number with $d$ as in~\eqref{eq:three}. Then $r$ is a root of $dx^3-(d-2)x-2n$, and $\lfloor(2n/d)^{1/3}\rfloor\in\{r-1,r\}$.
\end{lemma}
\begin{proof}
Multiplying the first equation of~\eqref{eq:three} by $r$ gives $2n=2Pqr=dr^3-(d-2)r$. Put $x=2n/d$; then $x-r^3=-(d-2)r/d\in(-r,r]$. If $x\ge r^3$ then $r\le x^{1/3}<(r^3+r)^{1/3}<r+1$; if $x<r^3$ then $x>r^3-r>(r-1)^3$ (as $3r^2-3r+1>r$ for $r\ge1$), so $r-1<x^{1/3}<r$. Thus $\lfloor x^{1/3}\rfloor$ is $r$ or $r-1$, and the same holds for $\lfloor\lfloor x\rfloor^{1/3}\rfloor$ since $\lfloor x\rfloor\le x<\lfloor x\rfloor+1$ changes the cube root by less than $1$ in the same direction.
\end{proof}

\begin{algorithm}[h]
\caption{Enumeration of four-prime Williams numbers (campaign 1: all $n\le B$; campaign 2: all $n$ with $p_1\le X$)}\label{alg}
\begin{algorithmic}[1]
\State \textbf{Input:} $B$ (campaign 1) or $X$ (campaign 2). Put $S=\lceil\sqrt{2p_1}\rceil$, $R_j=\max(2jS,8p_1)$; call $(m,j)$ \emph{admissible} if $|jm-T|\le R_j$ where $T=4P^2$.
\ForAll{primes $p_1\equiv1\ (12)$ with $p_1<B^{1/4}$ (campaign 1) or $p_1\le X$ (campaign 2)}
\ForAll{primes $p_2\equiv1\ (12)$, $p_1<p_2\le\sqrt{24p_1^3+3}$, and $p_2<(B/p_1)^{1/3}$ in campaign 1}
\State $P\gets p_1p_2$; \textbf{if} $p_1\mid p_2^2-1$ or $p_1\not\equiv p_2\pmod{\gcd(L_{p_1},L_{p_2})}$ \textbf{then skip} (Lemma~\ref{lem:crt})
\State compute $m_0,\Lambda$ by CRT; $m_{\max}\gets\min(4P^2-1,\lfloor B/P\rfloor)$ (no second term in campaign 2)
\State $M^*\gets\lfloor\sqrt{T}\sqrt{\Lambda}\rfloor$ (floating point; any value works, this balances the two walks)
\ForAll{$m\equiv m_0\ (\Lambda)$ with $p_2^2<m\le\min(M^*,m_{\max})$} \Comment{walk A}
\State keep $m$ iff $(m,\lfloor T/m\rfloor)$ or $(m,\lceil T/m\rceil)$ is admissible
\EndFor
\ForAll{integers $j$ from $\max\bigl(1,\lfloor(T-8p_1)/(m_{\max}+2S)\rfloor-1\bigr)$ to $\lfloor(T+8p_1)/(M^*-2S)\rfloor+1$} \Comment{walk B}
\State keep every $m\equiv m_0\ (\Lambda)$ with $\max(M^*+1,\lceil(T-R_j)/j\rceil)\le m\le\min(m_{\max},\lfloor(T+R_j)/j\rfloor)$, unless $(m,j-1)$ is admissible (reported once, under the least admissible $j$)
\EndFor
\ForAll{kept $m$} \Comment{survivors}
\State $n\gets Pm$; \textbf{if} $2^{n-1}\not\equiv1\pmod n$ \textbf{then discard} (in C with $128$-bit arithmetic when $P\le2^{40}$; otherwise deferred to the Python verifier)
\ForAll{$d=1,\dots,2p_1$} \State $x\gets\lfloor\lfloor2n/d\rfloor^{1/3}\rfloor$; \textbf{for} $r\in\{x-1,x,x+1\}$ with $dr^3-(d-2)r=2n$ and $r\mid m$: $q\gets m/r$; report $(p_1,p_2,q,r)$ if $p_2<q<r$, $q,r$ prime and all four conditions of Lemma~\ref{lem:quotient} hold
\EndFor
\EndFor
\EndFor
\EndFor
\end{algorithmic}
\end{algorithm}

\section{Completeness, validation and certification}\label{sec:complete}

\begin{proposition}[completeness]\label{prop:complete}
Let $n=p_1p_2qr$ be a four-prime Williams number with $n\le B$ (campaign 1) or $p_1\le X$ (campaign 2). Then Algorithm~\ref{alg} visits the pair $(p_1,p_2)$, keeps $m=qr$, and reports $n$.
\end{proposition}
\begin{proof}
By Theorem~\ref{thm:A}(a) both $p_1,p_2$ are primes $\equiv1\pmod{12}$. If $n\le B$ then $p_1^4<n\le B$ and $p_1p_2^3<n\le B$, so $p_1<B^{1/4}$ and $p_2<(B/p_1)^{1/3}$; in either campaign $p_2\le\sqrt{24p_1^3+1}$ by Theorem~\ref{thm:A}(c). Hence the pair is in the range scanned. The congruences of Lemma~\ref{lem:crt} hold for $m$, so the pair is not skipped and $m\equiv m_0\pmod\Lambda$; by Theorem~\ref{thm:A}(b), $p_2^2<m<4P^2$, and $m\le B/P$ when $n\le B$, so $m$ lies in the walked range. By Lemma~\ref{lem:window} the pair $(m,de)$ is admissible. If $m\le M^*$ then $\lfloor T/m\rfloor$ or $\lceil T/m\rceil$ is admissible and walk A keeps $m$. If $m>M^*$ then $j=de$ lies in the scanned range, because admissibility gives $jm\le T+2jS+8p_1$ and $jm\ge T-2jS-8p_1$, hence $j(m-2S)\le T+8p_1$ with $m-2S>M^*-2S$ and $j(m+2S)\ge T-8p_1$ with $m+2S\le m_{\max}+2S$; and $m$ lies in the exact window for the least admissible $j$, so walk B keeps it. Since $n$ is a Carmichael number, $2^{n-1}\equiv1\pmod n$. Finally $1\le d<2p_1$ by Theorem~\ref{thm:A}(b), and by Lemma~\ref{lem:cbrt} the integer cube root lands within one of $r$, so the loop over $d$ recovers $r$, $q=m/r$, and all checks pass.
\end{proof}

\paragraph{Validation.} The implementation was tested against brute force on a \emph{relaxed} problem that has solutions: replace $L_p$ by $(p^2-1)/\kappa$ with $\kappa\in\{24,120\}$ (Theorem~\ref{thm:A}(b) and Lemmas~\ref{lem:crt}--\ref{lem:cbrt}, which are all the algorithm uses, go through with $2$ replaced by $\kappa$; Theorem~\ref{thm:A}(c) and the Fermat filter are specific to $\kappa=2$ and are switched off in these runs). On the relaxed problem the enumerator finds exactly the solution set that a direct four-fold loop finds; for $\kappa=24$ the unique solution below $10^{12}$ is $13685=5\cdot7\cdot17\cdot23$, and the outputs of both programs are archived. Separate self-tests brute-force the two-prime sub-problem $(q,r)$ for small $(p_1,p_2)$ in four relaxed variants and check every inequality of Lemmas~\ref{lem:window}--\ref{lem:cbrt} on the solutions found, brute-force the residue class of Lemma~\ref{lem:crt} for $13{,}037$ pairs, and compare the $128$-bit Fermat routine with Python on random moduli. The two walks were checked to produce identical survivor lists on six instances, two of them relaxed instances with $346{,}409$ and $11{,}007{,}875$ survivors. These are regression tests of the implementation on instances with known answers, not an independent validation of the production campaigns, which rests on Proposition~\ref{prop:complete} and the certification below.

\paragraph{Certification.} Each campaign is processed in blocks of consecutive $p_1$ and logs one record per block (its $p_1$ range, number of pairs, survivors, Fermat survivors, solutions). Both campaigns were run twice, the second time with every window survivor written out, and a certification script asserts: every block identifier appears exactly once; block ranges are increasing, disjoint and contiguous over the admissible primes; per-block counts equal those of the first run; the distinct survivors, within each campaign, number exactly the logged total, block by block; every survivor re-tested in Python (base-2 Fermat test on $n=Pm$, then factorisation of $m$ and the four conditions of Lemma~\ref{lem:quotient} on the original tuple) yields the logged number of Fermat survivors and no Williams number. The records include hashes of the binary, the source and the logs, the compiler, host and source revision. An external reviewer repeated both certifications from the clean tagged checkout.

\section{Computational results for four prime factors}\label{sec:results}

\begin{theorem}\label{thm:B}
There is no Williams number with exactly four prime factors and $n\le10^{30}$, and there is none at all whose smallest prime factor is at most $500{,}000$. Consequently a four-prime Williams number, if one exists, has $p_1>500{,}000$ and $n>10^{30}$. (McIntosh and Dipra \cite{MD15} had $p_1>2\times10^4$ and $n>10^{24}$.)
\end{theorem}

\begin{proof}[Computation]
Table~\ref{tab:k4} summarises the two campaigns (Apple M4 Pro, 14 threads). In each, the single survivor of the Fermat filter is the Carmichael number $26921089=13\cdot37\cdot97\cdot577$, for which $38\nmid26921090$. The survivor totals are deduplicated within their respective campaigns. The second statement combines the second campaign with Corollary~\ref{cor:A}.
\end{proof}

\begin{table}[h]
\centering\small
\begin{tabular}{lrrrrr}
\toprule
campaign & pairs $(p_1,p_2)$ & window survivors & Fermat survivors & solutions & wall time\\
\midrule
all $n\le10^{30}$ & 230{,}612{,}582{,}193 & 1{,}904 & 1 & 0 & 87 min\\
all $n$, $p_1\le500{,}000$ & 86{,}517{,}004{,}140 & 1{,}185 & 1 & 0 & 146 min\\
\bottomrule
\end{tabular}
\caption{The four-prime enumeration (Theorem~\ref{thm:B}).}\label{tab:k4}
\end{table}

\section{No Williams number below \texorpdfstring{$10^{24}$}{10\^{}24}}\label{sec:bound}

\begin{theorem}\label{thm:bound}
No Williams number is smaller than $10^{24}$. Among the $308{,}279{,}939$ Carmichael numbers below $10^{24}$, exactly $14$ have four prime factors $p$ with $(p+1)\mid(n+1)$ (seven with $\omg(n)=9$, five with $\omg(n)=10$, two with $\omg(n)=11$), one of them
\[
15500309690554865758009=37\cdot103\cdot109\cdot127\cdot307\cdot523\cdot613\cdot1429\cdot2089
\]
(the primes $37,109,1429,2089$); none has five or more; $826$ have exactly three, $63{,}231$ exactly two and $5{,}807{,}346$ exactly one. Below $10^{22}$ the maximum is three, attained by $179$ numbers.
\end{theorem}

We know of no larger published bound that applies without restriction on the number of prime factors: Pinch's $10^{16}$ \cite{Pinch07} is restricted to an odd number of prime factors; McIntosh and Dipra \cite{MD15} give $10^{24}$ for four prime factors (indeed for Carmichael numbers with four prime factors whose largest prime $s$ satisfies $(s+1)\mid(N+1)$), note that Feitsma's list of base-2 pseudoprimes below $2^{64}$ contains no Williams number, and report a Pinch-style search for $4\le d\le13$ prime factors ``for various values of $X$ between $10^{20}$ and $10^{35}$'' without stating the bound reached for each $d$.

\subsection*{Input data} Carmichael numbers have been tabulated completely to $10^{21}$ by Pinch \cite{Pinch07,PinchTables}, to $10^{22}$ by Goutier \cite{Goutier} and to $10^{24}$ by Shallue and Webster \cite{SW24,SW25}, whose list is distributed by Stamm-Wilbrandt as a PARI/GP binary vector \cite{SW_bin}. A Williams number is in particular Carmichael, so it suffices to test every tabulated Carmichael number for the Lucas--Carmichael property. This is a factor-and-test pass over an external complete tabulation, not an independent tabulation of the Carmichael numbers; the tabulation's completeness is Shallue and Webster's result, which we check for consistency as described below.

\subsection*{Per-entry verification} Each $n$ was factored and verified by us rather than read from the tabulators' flags: Miller--Rabin square-root-of-unity splitting (for a Carmichael number, bases that are strong liars yield no split, so bases are tried until a witness produces a nontrivial square root of $1$; general factorisation is the fallback), then the factorisation was checked to be squarefree with product $n$ and to satisfy Korselt's criterion, and the number of primes $p$ with $(p+1)\mid(n+1)$ was recorded. Processing is line-based so that no entry can be skipped at a buffer boundary.

\subsection*{Integrity checks} The counts below $10^{22}$, $10^{23}$ and $10^{24}$ ($49{,}679{,}870$; $123{,}381{,}982$; $308{,}279{,}939$) equal the published ones \cite{SW25}; the list is strictly increasing and begins $561,1105,1729,2465,2821$; our per-$\omg$ tallies agree with independently published counts, $1{,}494{,}738$ three-prime Carmichael numbers below $10^{24}$ (OEIS A132195, \cite{Chick}) and $6{,}529{,}921$ eleven-prime ones (OEIS A299711); below $10^{22}$ all eleven per-$\omg$ counts agree with Goutier's tally, and the $179$ three-match numbers below $10^{22}$ agree with OEIS A329948 (Suteu; b-file by Eldar). The converse direction was also checked: none of the $129{,}125$ Lucas--Carmichael numbers below $10^{15}$ tabulated by Suteu (OEIS A006972) is Carmichael. Table~\ref{tab:omega} gives the distribution below $10^{24}$. The scan was run twice: the first run recorded the per-$\omg$ counts, the three-match count and one number attaining the maximum; the second (checkpointed) run recorded the full histogram of matching primes by $\omg(n)$ and is certified by a script that asserts a contiguous set of chunks, the total $308{,}279{,}939$, histogram sum equal to the total and zero verification failures; the two runs agree on every shared statistic.

\begin{table}[h]
\centering\small
\begin{tabular}{rr@{\qquad}rr@{\qquad}rr}
\toprule
$\omg(n)$ & count & $\omg(n)$ & count & $\omg(n)$ & count\\
\midrule
3 & 1{,}494{,}738 & 7 & 62{,}338{,}997 & 11 & 6{,}529{,}921\\
4 & 548{,}265 & 8 & 99{,}469{,}866 & 12 & 547{,}590\\
5 & 3{,}429{,}390 & 9 & 81{,}041{,}495 & 13 & 15{,}328\\
6 & 19{,}665{,}194 & 10 & 33{,}199{,}094 & 14 & 61\\
\bottomrule
\end{tabular}
\caption{Carmichael numbers below $10^{24}$ by number of prime factors, as recounted by our scan (total $308{,}279{,}939$).}\label{tab:omega}
\end{table}

\begin{corollary}\label{cor:apps}
\begin{enumerate}[label=\textup{(\alph*)}]
\item A counterexample to Agrawal's conjecture of the Lenstra--Pomerance type (a Williams number all of whose primes are $\equiv3\pmod{80}$) exceeds $10^{24}$.
\item Any counterexample to the Baillie--PSW test produced by Pomerance's construction \cite{Pomerance84} exceeds $10^{24}$, because that construction produces a Williams number.
\end{enumerate}
\end{corollary}
Statement (b) adds little to the known absence of Baillie--PSW counterexamples below $2^{64}\approx1.8\times10^{19}$ \cite{Gilchrist}, but it bounds a construction rather than a tested range.

\begin{remark}[Howe's order-$2$ Carmichael numbers]
By Howe's criterion \cite{Howe00}, a squarefree $n$ is a Carmichael number of order $2$ iff for each $p\mid n$, $n\equiv1$ or $n\equiv p\pmod{p^2-1}$. For a Williams number the first option is impossible ($n\equiv p\pmod{L_p}$ and $n\equiv1\pmod{p^2-1}$ would give $L_p\mid p-1$), so a Williams number is of order $2$ iff $n\equiv p\pmod{p^2-1}$ for all $p\mid n$, i.e. iff $(n-p)/L_p$ is even for every $p\mid n$. Tabulations of order-$2$ Carmichael numbers therefore bound one $2$-adic type of Williams number; we have not pursued this.
\end{remark}

\section{Discussion}\label{sec:disc}

\paragraph{A heuristic for four prime factors (speculative).} Condition on Theorem~\ref{thm:A}(a). For $p\equiv1\pmod{12}$ the condition $n/p\equiv1\pmod{L_p}$ is automatic modulo $12$; treat the remaining condition as an event of probability $24/(p^2-1)\approx24/p^2$, independently for the four primes (the four events are in fact arithmetically dependent, and no error term is claimed). Summing $\prod_i24/p_i^2$ over quadruples of primes $\equiv1\pmod{12}$ with $p_2^2\le24p_1^3+1$ and $p_2<q<r<2p_1p_2$ gives about $6.3\times10^{-7}$ (script \texttt{agrawal/k4\_heuristic.py}), of which about $96\%$ comes from $p_1=13$; the contribution of $p_1>X$ decays like $X^{-4}$ up to logarithms. On this model four-prime Williams numbers are not expected, and Theorem~\ref{thm:B} is consistent with the prediction.

\paragraph{Open problems.}
\begin{enumerate}[leftmargin=2em]
\item Is $k=4$ impossible? By Theorem~\ref{thm:A}, fixing $p_1$ leaves a finite search over $(p_2,q,r)$; the integer $\delta=des-8p_1^3=8p_1^3(1/\Pi-1)$ satisfies $\delta\ge1$ \cite[proof of Thm.~2]{MD15} and ties the sizes of $p_2,q,r$ to $p_1$ but does not by itself exclude them.
\item The $k=5$ case, where both classes modulo $4$ are allowed and the analogue of~\eqref{eq:des} has four factors.
\item Closing the $m=4$ polynomial-family stratum with one repeated root (Appendix~\ref{sec:obstr}).
\end{enumerate}

\paragraph{Reproducibility.}\label{sec:repro}
All code is in the repository \url{https://github.com/razv2025/carmichael-lucas-carmichael}, tagged for this version. Inventory of the computations and their verification status: (i) the two four-prime campaigns were run twice, certified as described in \S\ref{sec:complete}, and the certifications were repeated by an external reviewer from the tagged source; (ii) the $10^{24}$ scan was run twice (the second run with the full match histogram and a certification record), with the integrity checks of \S\ref{sec:bound} against externally published counts, and its input is identified by SHA-256; (iii) the relaxed-problem validation to $10^{12}$ and the self-tests are archived and re-runnable (\texttt{make test}, under a minute); (iv) the polynomial-family searches of Appendix~\ref{sec:obstr} are archived with their command lines. The four-prime enumeration (\texttt{clc/k4/}) consists of the C program, its brute-force reference for the relaxed problem, the Python re-verifier, the per-block logs and survivor dumps of both campaigns, the certification records and \texttt{make test}; the $10^{24}$ scan (\texttt{clc/}) consists of the download--decode--verify--scan pipeline, the completeness report and the result files. The code was developed with AI assistance; see the declaration below.

\appendix

\section{Polynomial families}\label{sec:obstr}

Carmichael numbers are constructed in two ways: by polynomial families in the manner of Chernick \cite{Chernick39}, $n(k)=(6k+1)(12k+1)(18k+1)$, where the divisibilities are polynomial identities; and by subset products in the manner of Erd\H{o}s and Alford--Granville--Pomerance \cite{Erdos56,AGP94,LN96}. This appendix shows that the first construction cannot be made two-sided in its simplest forms. \emph{The theorems concern families whose divisibilities hold identically in the parameter $k$; they say nothing about families that are Williams numbers only for $k$ in some residue classes, nor about sporadic values.} The second construction is discussed, with measurements, in the supplementary note \cite{Supp}.

Consider $n(k)=\prod_{i=1}^m(a_ik+b_i)$ with integers $a_i>0$, the linear factors pairwise non-proportional, and require $(a_ik+b_i-1)\mid(n(k)-1)$ and $(a_ik+b_i+1)\mid(n(k)+1)$ in $\mathbb Q[k]$. Since a linear polynomial divides $f$ iff $f$ vanishes at its root, this is $n\bigl((1-b_i)/a_i\bigr)=1$ and $n\bigl((-1-b_i)/a_i\bigr)=-1$. With $D_{ij}=a_ib_j-a_jb_i$ one computes $a_j\frac{1-b_i}{a_i}+b_j=(a_j+D_{ij})/a_i$, so the conditions read, for every $i$,
\begin{equation}\label{eq:fam}
\prod_{j\ne i}(a_j+D_{ij})=a_i^{m-1}=\prod_{j\ne i}(D_{ij}-a_j).
\end{equation}
If $g=\gcd(a_1,\dots,a_m)>1$, write $a_i=ga_i'$; then $D_{ij}=gD'_{ij}$ and~\eqref{eq:fam} for $(a,b)$ is $g^{m-1}$ times~\eqref{eq:fam} for $(a',b)$, so every family reduces to one with $\gcd(a_i)=1$ (the same numbers on the sub-progression $k\in g\mathbb Z$); the searches below therefore assume $\gcd(a_i)=1$.

\begin{theorem}[three factors]\label{thm:fam3}
No such family exists for $m=3$.
\end{theorem}
\begin{proof}
Subtracting the two products in~\eqref{eq:fam} for $i=1,2,3$ gives $a_2D_{13}+a_3D_{12}=0$, $a_1D_{23}+a_3D_{21}=0$, $a_1D_{32}+a_2D_{31}=0$, i.e. $a_2D_{13}=-a_3D_{12}$, $a_1D_{23}=a_3D_{12}$, $a_1D_{23}=-a_2D_{13}$. The identity $a_1D_{23}-a_2D_{13}+a_3D_{12}=0$ (it is $a\wedge b\wedge a=0$) becomes $3a_3D_{12}=0$, so $D_{12}=0$ and then $D_{13}=D_{23}=0$: $b$ is proportional to $a$ and the factors are proportional, contrary to hypothesis.
\end{proof}

\begin{theorem}[uniform offsets]\label{thm:fampm1}
No such family exists with $b_i=-1$ for all $i$, nor with $b_i=+1$ for all $i$, for any $m\ge2$.
\end{theorem}
\begin{proof}
If all $b_i=-1$ then $D_{ij}=a_j-a_i$, the second product in~\eqref{eq:fam} is $(-a_i)^{m-1}$, and the first becomes $\prod_{j\ne i}(2a_j-a_i)=a_i^{m-1}$. The $a_i$ are distinct (non-proportionality). Take $i$ with $a_i$ maximal: for $j\ne i$, $0<a_j<a_i$ gives $|2a_j-a_i|<a_i$, so the product has absolute value $<a_i^{m-1}$. The case $b_i=+1$ is the mirror image with $\prod_{j\ne i}(a_i-2a_j)$.
\end{proof}

\begin{theorem}[generic position]\label{thm:famgen}
Put $s_i=(b_i-1)/a_i$ and $t_i=(b_i+1)/a_i$. If the $s_i$ are pairwise distinct and the $t_i$ are pairwise distinct, no such family exists, for any $m\ge2$.
\end{theorem}
\begin{proof}
Let $G(\tau)=n(-\tau)$, of degree $m$ with leading coefficient $c=(-1)^m\prod a_i$. The hypotheses say $G(s_i)=1$ and $G(t_i)=-1$. With $m$ distinct roots each, $G-1=c\prod(\tau-s_i)$ and $G+1=c\prod(\tau-t_i)$. Subtracting, $2=c\bigl[\prod(\tau-t_i)-\prod(\tau-s_i)\bigr]$ identically in $\tau$. The coefficient of $\tau^{m-1}$ on the right is $-c\sum(t_i-s_i)=-c\sum 2/a_i\ne0$, while on the left it is $0$.
\end{proof}

\begin{corollary}
Every such family has two indices $i\ne j$ with $p_j-1=\frac{a_j}{a_i}(p_i-1)$ or $p_j+1=\frac{a_j}{a_i}(p_i+1)$ as polynomial identities.
\end{corollary}

The strata left open by Theorem~\ref{thm:famgen} were searched exhaustively using the fact that each factor in~\eqref{eq:fam} divides $a_i^{m-1}$: no family exists with $m=4$ and all $a_i\le500$, or with $m=5$ and $a_i\le150$; Theorem~\ref{thm:fampm1} was also confirmed by search for $m=5,7,9,11$ up to $a_i\le20000,5000,3000,1500$. The outputs of these runs, with their command lines, are archived in \texttt{agrawal/results/}. Theorem~\ref{thm:fam3} is the parametric shadow of Proposition~\ref{prop:k3}.

\section*{Declaration of generative AI and AI-assisted technologies}
An AI system (Claude, Anthropic) was used, under the author's direction, for: drafting and checking the proofs; writing the code (the enumerator, the scanners, the certification and test scripts); drafting this text; literature search; and adversarial checking of the proofs and of the numerical claims. The author reconstructed and checked every proof and every claim, reviewed the code, and accepts full responsibility for the content. AI checks are not a substitute for human review; the human reviews received are acknowledged below.

\section*{Acknowledgements}
We thank the authors of the tabulations this work rests on: R.~G.~E.~Pinch, C.~Goutier, A.~Shallue and J.~Webster, H.~Stamm-Wilbrandt (for making the $10^{24}$ list available), D.~Suteu, and the OEIS contributors whose independent counts served as integrity checks; and the reviewers who read earlier drafts, pointed us to the Williams-number literature, and repeated the certifications.

\begin{thebibliography}{99}
\bibitem{AG13} I.~Al-Rasasi, N.~Ghanmi, On Williams numbers with three prime factors, \emph{Missouri J. Math. Sci.} 25 (2013), 134--152. \url{https://doi.org/10.35834/mjms/1384266199}
\bibitem{AKS} M.~Agrawal, N.~Kayal, N.~Saxena, PRIMES is in P, \emph{Ann. of Math.} 160 (2004), 781--793. \url{https://doi.org/10.4007/annals.2004.160.781}
\bibitem{AGP94} W.~R.~Alford, A.~Granville, C.~Pomerance, There are infinitely many Carmichael numbers, \emph{Ann. of Math.} 139 (1994), 703--722. \url{https://doi.org/10.2307/2118576}
\bibitem{BEP10} K.~Bouall\`egue, O.~Echi, R.~G.~E.~Pinch, Korselt numbers and sets, \emph{Int. J. Number Theory} 6 (2010), 257--269. \url{https://doi.org/10.1142/S1793042110002922}
\bibitem{BFW21} R.~Baillie, A.~Fiori, S.~S.~Wagstaff Jr., Strengthening the Baillie--PSW primality test, \emph{Math. Comp.} 90 (2021), 1931--1955. \url{https://doi.org/10.1090/mcom/3616}
\bibitem{ChenGreene03} Z.~Chen, J.~Greene, Some comments on Baillie--PSW pseudoprimes, \emph{Fibonacci Quart.} 41 (2003), 334--344.
\bibitem{Chernick39} J.~Chernick, On Fermat's simple theorem, \emph{Bull. Amer. Math. Soc.} 45 (1939), 269--274. \url{https://doi.org/10.1090/S0002-9904-1939-06953-X}
\bibitem{Chick} J.~M.~Chick, Carmichael number variable relations: three-prime Carmichael numbers up to $10^{24}$, \texttt{arXiv:0711.2915} (2007); counts in OEIS A132195.
\bibitem{Echi07} O.~Echi, Williams numbers, \emph{C. R. Math. Rep. Acad. Sci. Canada} 29 (2007), 41--47.
\bibitem{Erdos56} P.~Erd\H{o}s, On pseudoprimes and Carmichael numbers, \emph{Publ. Math. Debrecen} 4 (1956), 201--206.
\bibitem{Gilchrist} J.~Gilchrist, Pseudoprime enumeration with probabilistic primality tests, web page (2013), reporting the 2009 verification, from J.~Feitsma's table of base-2 pseudoprimes below $2^{64}$, that no Baillie--PSW pseudoprime exists below $2^{64}$; see also T.~R.~Nicely's Baillie--PSW page.
\bibitem{Goutier} C.~Goutier, Carmichael numbers up to $10^{22}$ (2024), data and readme linked from OEIS A002997 and A055553.
\bibitem{Grantham} J.~Grantham, Frobenius pseudoprimes, \emph{Math. Comp.} 70 (2001), 873--891; and the \$620 problem page.
\bibitem{Howe00} E.~W.~Howe, Higher-order Carmichael numbers, \emph{Math. Comp.} 69 (2000), 1711--1719. \url{https://doi.org/10.1090/S0025-5718-00-01225-4}
\bibitem{LN96} G.~L\"oh, W.~Niebuhr, A new algorithm for constructing large Carmichael numbers, \emph{Math. Comp.} 65 (1996), 823--836. \url{https://doi.org/10.1090/S0025-5718-96-00692-8}
\bibitem{LP} H.~W.~Lenstra Jr., C.~Pomerance, Remarks on Agrawal's conjecture, in: \emph{Future directions in algorithmic number theory}, AIM workshop notes (2003), pp.~30--32; archived copy of \url{http://aimath.org/WWN/primesinp/primesinp.pdf}.
\bibitem{McIntosh14} R.~J.~McIntosh, Carmichael numbers with $(p+1)\mid(n-1)$, \emph{Integers} 14 (2014), \#A59.
\bibitem{MD15} R.~J.~McIntosh, M.~Dipra, Carmichael numbers with $p+1\mid n+1$, \emph{J. Number Theory} 147 (2015), 81--91. \url{https://doi.org/10.1016/j.jnt.2014.06.005}
\bibitem{Phatak20} D.~Phatak, Bounds on the range(s) of prime divisors of a class of Baillie--PSW pseudo-primes, \texttt{viXra:2007.0102} (2020), unrefereed.
\bibitem{Pinch07} R.~G.~E.~Pinch, Absolute quadratic pseudoprimes, in: \emph{Proc. Conf. on Algorithmic Number Theory}, Turku 2007, TUCS General Publication 46, pp.~113--128.
\bibitem{PinchTables} R.~G.~E.~Pinch, The Carmichael numbers up to $10^{21}$, \emph{Proc. Conf. on Algorithmic Number Theory}, Turku 2007; tables at \url{http://www.s369624816.websitehome.co.uk/rgep/carpsp.html}
\bibitem{Pomerance84} C.~Pomerance, Are there counter-examples to the Baillie--PSW primality test?, in: \emph{Dopo le parole} (H.~W.~Lenstra Jr., J.~K.~Lenstra, P.~van Emde Boas, eds.), Amsterdam 1984. \url{https://www.pseudoprime.com/dopo.pdf}
\bibitem{SW24} A.~Shallue, J.~Webster, Advances in tabulating Carmichael numbers, \emph{Res. Number Theory} 11 (2025), art.~8. \url{https://doi.org/10.1007/s40993-024-00598-3}
\bibitem{SW25} A.~Shallue, J.~Webster, Algorithms for Carmichael numbers, \texttt{arXiv:2506.09903} (2025).
\bibitem{SW_bin} H.~Stamm-Wilbrandt, \texttt{carm10e24.bin}, PARI/GP binary vector of all Carmichael numbers below $10^{24}$ (7 March 2026). \url{https://stamm-wilbrandt.de/carm10e24.bin}
\bibitem{Supp} R.~Vakil, Supplementary note: two-sided subset products and an onset heuristic for Williams numbers; file \texttt{paper/supplement.pdf} in the repository.
\bibitem{Williams77} H.~C.~Williams, On numbers analogous to the Carmichael numbers, \emph{Canad. Math. Bull.} 20 (1977), 133--143. \url{https://doi.org/10.4153/CMB-1977-025-9}
\bibitem{Zhang15} Z.~Zhang, Estimating the counts of Carmichael and Williams numbers with small multiple seeds, \emph{Math. Comp.} 84 (2015), 309--337. \url{https://doi.org/10.1090/S0025-5718-2014-02837-8}
\end{thebibliography}
\end{document}
